\documentclass{article}
\usepackage[T1]{fontenc}
\usepackage{lmodern}
\usepackage{graphicx}
\usepackage{amsmath,amssymb}
\usepackage[a4paper,margin=2cm]{geometry}
\usepackage[absolute,overlay]{textpos}
\setlength{\TPHorizModule}{1mm}
\setlength{\TPVertModule}{1mm}
\usepackage[indonesian]{babel}
\usepackage{hyperref}
\setlength{\emergencystretch}{2em}
% unit: Halaman 5

\begin{document}

% === Halaman 5 (llm) ===

\section*{Akar Primitif dan Logaritma Diskrit}

\begin{itemize}
    \item Jika $n$ adalah bilangan bulat, maka $a$ disebut \textbf{akar primitif} dari $n$ jika perpangkatan
    \[ a, a^2, ..., a^{\phi(n)} \quad (\text{dalam sebuah modulus } n) \]
    menghasilkan nilai yang berbeda dan semuanya relatif prima dengan $n$.
    
    \item Secara khusus, jika $p$ adalah bilangan prima, maka $a$ disebut akar primitif dari $p$ jika perpangkatan
    \[ a, a^2, ..., a^{p-1} \quad (\text{dalam modulus } p) \]
    menghasilkan nilai-nilai yang berbeda \\
    (ingatlah dari fungsi \textit{toitient} Euler, bahwa jika $p$ prima maka $\phi(p) = p - 1$)
\end{itemize}

\end{document}
